\section{Introduction}

A contraction remembers an exceptional cycle in several ways.
Its local quotient describes the cycle and its incidence with
neighboring fibres. Its embedding in the compact space determines
integral intersection and attachment data. A degenerating Kähler
metric measures how that same cycle disappears in distance.
The relation between these descriptions is especially concrete
for moduli of bundles on a K3 surface: stable factors determine
a quiver, while universal families identify the quiver's cycles
inside integral cohomology.

This paper joins these descriptions for compact contractions
without bubbling. The central example has a moving projective
line and a deeper Grassmannian fibre \(\Gr(2,4)\).
The line specializes to a Schubert line. The two middle
Schubert classes then reveal information absent from divisor
restrictions: they prove integral primitivity and calculate an
attachment modulo six. On these same cycles, the actual
harmonic metric has exact periods and explicit ambient diameter
asymptotics. A general comparison theorem controls the graph
of the entire compact contraction.

\subsection{The compact Grassmannian example}
The Mukai pairing is
\[
 (r,c,s)\cdot(r',c',s')=c\cdot c'-rs'-r's .
\]
All cohomology and homology in the following statement have
integral coefficients. The harmonic metric is the real,
unnormalized Dolbeault \(L^2\) metric on moduli.

For a universal topological \(K\)-class \(\mathcal E\) on
\(X\times M\), write
\[
 I_z=p_!(z^\vee\otimes\mathcal E),\qquad z\in K^0_{\rm top}(X),
\]
where \(p:X\times M\to M\) and \(p_!\) is the topological
\(K\)-theory pushforward. Thus \(z\) ranges over all integral
topological classes, not only algebraic ones. The marked index-Chern
classes on a complement are the restrictions of \(c_k(I_z)\).
The normalization used below tensors \(\mathcal E\) by the inverse
determinant of a specified rank-one index, as constructed in
Section~\ref{sec:schubert}.

\begin{theorem}\label{thm:main}
There is a projective K3 surface with
\[
 \NS(X)=\Z H\oplus\Z D,\qquad H^2=2,\quad H\cdot D=0,\quad
 D^2=-44,\quad H\text{ ample},
\]
and a unique \(H\)-stable spherical bundle \(A\) of vector
\(a=(3,D,-7)\). Put \(b=(1,0,1)\), \(w=a+b\),
and \(v=2a+b=(7,2D,-13)\).
For sufficiently small generic \(t\ne0\), the entire moduli
space \(M_\pm\) of vector \(v\), for polarization \(H+tD\),
is a smooth compact eightfold of \(K3^{[4]}\) type consisting
of slope-stable bundles. The two chamber spaces have projective
symplectic contractions \(\pi_\pm:M_\pm\to Z\) to a common
normal target with the following properties.

\begin{enumerate}
\item The reduced singular locus is a normal sixfold \(T\)
with one singular point \(o\). Its smooth part parametrizes
classes \(A\oplus F\), where \(F\) is an \(H\)-stable bundle
of vector \(w\); \(o\) represents \(2A\oplus\OO_X\).
The fibres of \(\pi_\pm\) are points, projective lines, and
\(G=\Gr(2,4)\), respectively.
The auxiliary generic \(w\)-moduli resolves \(T\), with
exceptional fibre \(\PP^3\).

\item A moving projective line specializes, in a proper reduced
family, to a minimal Schubert line \(\ell\subset G\).
Every homology inclusion \(H_k(G)\to H_k(M_\pm)\) is primitive,
and \([G]^2=6\). In the positive chamber, let
\(h=c_1(\mathcal U)\), \(p=c_2(\mathcal U)\) for the positive
tautological quotient on \(G\). Normalized universal indices
satisfy
\[
 c(I_z)|_G=(1+h+p)^{\chi(A,z)},\qquad
 \int_Gh^4=2,\quad\int_Gh^2p=\int_Gp^2=1 .
\]
For \(V=M_+\setminus G\), these formulas calculate
\(\kappa:H^8(V)\to\Z/6\) on all marked index-Chern generators,
and
\[
 H^8(M_+)\simeq
 \{(u,N)\in H^8(V)\oplus\Z:\ N\equiv\kappa(u)\pmod6\}.
\]
The marked complement ring remains input to this calculation.

\item Put \(\kappa_t=88\pi|t|/7\).
The minimal line and the moving fibre both have harmonic area
\(\pi\kappa_t\). Either Schubert plane has volume
\((\pi\kappa_t)^2/2\), and \(G\) has volume
\((\pi\kappa_t)^4/12\). Their ambient diameters satisfy
\[
 \diam G=\frac{\pi}{\sqrt2}\sqrt{\kappa_t}(1+O(|t|)),\qquad
 \diam\PP^2=\frac{\pi}{2}\sqrt{\kappa_t}(1+O(|t|)).
\]
The ambient spaces, rescaled by \(\kappa_t^{-1/2}\) and based
on \(G\), converge in complete pointed Gromov--Hausdorff
topology to the unit-level cotangent Grassmannian quotient.
The harmonic metric on \(Z\setminus T\) completes to \(Z\);
the graphs of \(\pi_\pm\) satisfy
Theorem~\ref{thm:quantitative}'s distortion bound and give
compact Gromov--Hausdorff convergence to that completion.
\end{enumerate}
There is no additional non-locally-free or positive-charge
boundary in this compact problem.
\end{theorem}

The local incidence geometry of the Grassmannian is classical.
The point of the theorem is its realization as the entire deep
fibre of this bundle-only compact family, together with its
integral attachment and actual harmonic geometry.
The equal-dimensional incidence resolutions for the two signs
are isomorphic over the affine cone; a nontrivial
Grassmannian flop is not part of the statement.

\subsection{The general reconstruction and metric comparison}
The general hypotheses, stated in
Section~\ref{subsec:hypotheses}, require compact chamber
spaces with a common central slope target, absence of every
ideal boundary, and full convergent quiver charts with a
module-simple point in each complex zero fibre.
The local charts include repeated factors, tangent directions,
inverse solution coverage and the actual classifying maps.

Two kinds of global information enter. For all-degree additive
reconstruction, retain the geometric regular locus and the
local classifying families, with their isomorphism loci.
These determine the actual multiple overlaps and monodromy.
Their integral cochain diagram computes the chamber cohomology
by Proposition~\ref{prop:descent}.
For a compact core \(E\), Theorem~\ref{thm:attachment}
instead uses its intersection map, calculated universal-index
restrictions, and the marked ring of \(M\setminus E\).
Under the stated middle-degree vanishing hypotheses it computes
the boundary homomorphism and an integral fibre product.
These are two useful levels of reconstruction; neither
recovers an unknown global complement from singularity types.

The metric result is quantitative in finite-model distances.
For a sequence of prescribed accuracies we construct finite
covers, positive short-distance scales \(\ell_j\), relative
errors \(\alpha_j\to0\), and variational discrepancies
\(\beta_j(\theta)\to0\) for each fixed \(j\).
If \(R_\theta\) is the graph of the actual contraction, then
\[
 \dis R_\theta\le
 \inf_{j:\theta\in\mathcal P_j}
 \left[\left(\alpha_j+\frac{\beta_j(\theta)}{\ell_j}\right)D_*
                +\beta_j(\theta)\right],
\]
where \(D_*\) is constructed from a fixed central chart cover.
Definitions \eqref{eq:finite-b}--\eqref{eq:global-parameters}
make all these quantities explicit.
The estimate is uniform through stabilizer changes and
arbitrary adjacent-chamber approaches. It is a variational
modulus, rather than a universal power-law rate.

Three points in the proof deserve emphasis.
First, no-charge compactness and local classifying compatibility
put \emph{whole} fibres in the analytic charts.
Second, paths through lower strata can be approximated in
length by paths in the principal stratum; density alone would
not identify the harmonic completion.
Third, first-exit estimates compare both ambient distance
infima, including paths which might leave a local model.
A short-distance correspondence lemma then assembles these
local estimates without using a qualitative global convergence
theorem.

\subsection{Relation to earlier work}
Existence and deformation type of the K3 moduli spaces come
from Yoshioka \cite{yoshioka99,yoshioka01}.
The slope compactifications and the surface comparison with
gauge theory are due to Greb--Toma and
Greb--Sibley--Toma--Wentworth \cite{grebtoma,gstw}.
Quiver descriptions and wall geometry are developed in
Arbarello--Saccà and Bayer--Macrì
\cite{arbarellosacca,bayermacri}.
Grassmannian geometry in K3 moduli also appears in Markman's
Brill--Noether duality \cite{markmanbn}; the present compact
realization is proved from its own rank and charge classification.

The integral generation theorem is Markman's
\cite{markman}, and cochain descent is classical.
Our attachment theorem isolates the elementary relative
argument that turns \emph{calculated} index restrictions into
a boundary homomorphism respecting every integral relation.
The mod-six application uses both Schubert directions; it
does not claim a new presentation of Grassmannian cohomology.
Kaplan--Schedler \cite{kaplanschedler} study compatibility
and monodromy in gluing crepant resolutions.
Here existing compact chamber spaces and their geometric
classifying data are hypotheses; our cochain construction
retains, rather than solves for, their monodromy.

The qualitative formality theorem of
Bandiera--Manetti--Meazzini \cite{bmm} already applies to the
common-slope direct sums used here because their automorphism
groups are reductive. The controlled convergent symplectic
coordinates are imported from \cite{wuebben17}.
Relative finite balancing, complete quotient continuity,
normal slices and the harmonic comparison are imported from
\cite{wuebben14,wuebben15}; their hypotheses and normalizations
are restated in Section~\ref{sec:data}.
The complete cotangent model metric is classical
\cite{biquardgauduchon}.
Our metric contribution is the whole-fibre and principal-path
argument, the explicit global graph bound, and its application
to these identified compact cycles.
Quantitative comparison results such as
Lakzian--Sormani \cite{lakziansormani,lakziansormanier}
place ambient distance discrepancies among their inputs;
we derive the discrepancies here from finite quotient models.
The established qualitative Ricci-flat degeneration theory
\cite{tosatti,song} provides context, but is not an input to
Theorem~\ref{thm:quantitative}.

Sections~\ref{sec:integral}--\ref{sec:global} prove the general
statements. Sections~\ref{sec:construction}--\ref{sec:cycle-metrics}
prove Theorem~\ref{thm:main} on one compact family.
Section~\ref{sec:examples} recovers the fourfold applications
and uses the same two Schubert classes to exhibit the necessity
of retaining monodromy.
